% CPSY 1291 -- Assignment 1 call sheet
% One page: the signatures and idioms Recitations 1-2 taught, as a lookup card.
% Compile:  lualatex callsheet-A1.tex
\documentclass[10pt]{article}
\usepackage{cpsy1291-handout}
\usepackage{multicol}
\geometry{margin=0.55in}
\pagestyle{empty}
\setlength{\columnsep}{24pt}
\setlength{\parskip}{2pt}

\lstset{basicstyle=\ttfamily\footnotesize, aboveskip=3pt, belowskip=2pt, framesep=4pt}

% Compact block header and "what it returns / what to remember" line
\newcommand{\block}[1]{\vspace{5pt}{\sffamily\bfseries\color{brown}#1}\par\vspace{0pt}}
\newcommand{\ret}[1]{{\footnotesize\color{ink}#1\par}\vspace{1pt}}

\begin{document}

\begin{center}
{\sffamily\small\color{accent}\MakeUppercase{CPSY 1291 \ $\cdot$\ Assignment 1 call sheet}}\\[3pt]
{\sffamily\LARGE\bfseries\color{brown} Signatures and idioms}\\[2pt]
{\color{lightgray}\small Fall 2026 \quad$\cdot$\quad everything here was taught in Recitations 1--2 \quad$\cdot$\quad keep it next to the notebook}
\end{center}
\vspace{1pt}\hrule\vspace{2pt}

\begin{multicols}{2}

\block{Shapes, \texttt{axis=}, \texttt{keepdims}}
\begin{lstlisting}
X.shape                    # (120, 4096) stimuli x features
X.mean(axis=0).shape       # (4096,)  the named axis disappears
X.mean(axis=1, keepdims=True).shape   # (120, 1)
X - X.mean(axis=1)         # ValueError: (120,4096) (120,)
\end{lstlisting}
\ret{Compare shapes \emph{from the right}: each pair equal, or one is 1, or you get the error above. \texttt{keepdims=True} keeps a size-1 axis so subtraction lines up; \texttt{v[:, None]} is the same idea by hand.}

\block{Boolean masks}
\begin{lstlisting}
mask = (taxon == 2)     # one boolean per stimulus
X[mask]                 # just those rows
mask.sum(); mask.mean() # how many; what FRACTION (a rate)
np.isin(cat, [0, 2])    # membership in a set of labels
\end{lstlisting}
\ret{Element-wise logic on arrays: \texttt{\&} and \ $\cdot$\ \texttt{|} or \ $\cdot$\ \texttt{\textasciitilde} not --- \emph{always with parentheses}: \texttt{(taxon == 2) | (taxon == 5)}. Writing \texttt{and}/\texttt{or} between arrays raises ``truth value \dots\ is ambiguous''.}

\block{The upper triangle, and reordering a matrix}
\begin{lstlisting}
iu = np.triu_indices(120, k=1)  # a PAIR of index arrays
D[iu]                    # 7,140 values: every pair once
order = np.argsort(taxon)
D[np.ix_(order, order)]  # rows AND columns reordered
D[order, order]          # NOT that: 120 diagonal entries
\end{lstlisting}
\ret{\texttt{iu} walks the upper triangle in a \emph{fixed pair order}, so \texttt{D1[iu]} and \texttt{D2[iu]} describe the same pair at every position --- the whole reason two such vectors can be correlated. \texttt{np.ix\_} builds the open grid; \texttt{D[order, order]} silently applies the paired rule instead.}

\block{Fancy indexing and \texttt{argsort}}
\begin{lstlisting}
D[a, c] > D[a, b] + D[b, c] + 1e-9  # one boolean per triple
np.argsort(v)[:5]        # indices of the 5 SMALLEST
np.argsort(v)[::-1][:5]  # indices of the 5 LARGEST
\end{lstlisting}
\ret{Two index arrays pull one element \emph{per pair} --- 200{,}000 triples in one line, no loop (compute the distance matrix \emph{once}, then index into it). \texttt{argsort} returns \emph{indices}, not values --- and a stimulus is always maximally similar to itself, so exclude it from its own neighbor list.}

\block{Comparing two RDMs}
\begin{lstlisting}
rho = spearmanr(D1[iu], D2[iu]).statistic
\end{lstlisting}
\ret{\texttt{.statistic} is where the number lives (\texttt{.pvalue} exists, but the pairs are not independent --- do not lean on it). Spearman = Pearson on \emph{ranks}: only the ordering of dissimilarities is meaningful. Never feed the whole matrix: that counts every pair twice and adds $n$ zeros.}

\block{Condensed vs square}
\begin{lstlisting}
d = pdist(X, metric='euclidean')  # condensed (n*(n-1)/2,)
D = squareform(d)                 # square (n, n), and back
\end{lstlisting}
\ret{\texttt{imshow} and \texttt{MDS} want \emph{square}; \texttt{linkage} wants \emph{condensed}: \texttt{linkage(squareform(D\_human, checks=False), method='average')}.}

\columnbreak

\block{Fitting a curve, and the straight line}
\begin{lstlisting}
popt, _ = curve_fit(f_exp, x, y, p0=[10, 1], maxfev=40000)
y_hat   = f_exp(x, *popt)  # * unpacks the parameters back in
coef  = np.polyfit(x, y, 1)  # slope and intercept
y_hat = np.polyval(coef, x)  # the line at every x
\end{lstlisting}
\ret{\texttt{p0}: start near the right scale --- similarities run 0--10, so $a \approx 10$. \texttt{maxfev}: raise it, or hard fits die with ``Optimal parameters not found''. Wrap in \texttt{try/except RuntimeError} where a fit is \emph{allowed} to fail (2j says so --- that failure is a finding, not a bug).}

\block{Goodness of fit}
\[
R^2 = 1 - \frac{\sum_i (y_i - \hat{y}_i)^2}{\sum_i (y_i - \bar{y})^2}
\]
\ret{$1$ = perfect \ $\cdot$\ $0$ = no better than predicting the mean \ $\cdot$\ negative = worse than that (a real outcome, not a bug). It is a \emph{ratio}: two $R^2$ values are comparable only against the same target --- binning first raises it with the curve unchanged.}

\block{The scikit-learn pattern}
\begin{lstlisting}
est = PCA(n_components=50)
Y   = est.fit_transform(X)     # items x features in
est.explained_variance_ratio_  # exists only after .fit
cum = np.cumsum(est.explained_variance_ratio_)
np.argmax(cum >= 0.85)  # first True: components to 85%
\end{lstlisting}
\ret{Settings go in the constructor, data goes to \texttt{.fit}, learned results grow a \emph{trailing underscore} (not a typo). \texttt{MDS} takes a dissimilarity matrix: \texttt{metric='precomputed'}. Random methods differ between runs unless you set \texttt{random\_state} --- reproducible is not the same as meaningful.}

\block{Per-group anything (groupby without groupby)}
\begin{lstlisting}
means = np.stack([F[obj == k].mean(axis=0)
                  for k in range(51)])       # (51, 4096)
\end{lstlisting}
\ret{Four moves: mask $\rightarrow$ those rows $\rightarrow$ collapse \texttt{axis=0} $\rightarrow$ \texttt{np.stack}. Every ``for each object / category / cluster'' sentence is this template with a different word in step 3.}

\block{NaNs --- pick one, and say which}
\begin{lstlisting}
np.isnan(A).sum()  # count them before trusting anything
IT = it_raw[:, ~np.isnan(it_raw).any(axis=0)]     # DROP
IT = np.where(np.isnan(it_raw),                   # IMPUTE
              np.nanmean(it_raw, axis=0), it_raw)
\end{lstlisting}
\ret{Both are defensible; the notebook only demands you \emph{say which}. Ordering rule: check for zero-variance units \emph{before} z-scoring --- dividing by a std of 0 makes NaNs, and one NaN poisons everything downstream. Related clip: \texttt{np.sqrt(np.maximum(D2, 0))} --- floating point leaves tiny negatives where the true value is 0.}

\block{\texttt{silhouette\_score} --- what goes in}
\begin{lstlisting}
silhouette_score(Y[:, :2], (cat == c).astype(int))
\end{lstlisting}
\ret{First argument: the 2-D \emph{coordinates} in the PC1--PC2 plane, \emph{not} the RDM --- it computes its own distances. Second: a binary labeling, this category against everything else. Near 1 = sits apart \ $\cdot$\ near 0 = on a boundary \ $\cdot$\ negative = closer to another group than its own. No relation to 3f's silhouette \emph{masks} --- the name is a coincidence.}

\end{multicols}

\vspace{1pt}\hrule\vspace{4pt}
{\color{lightgray}\footnotesize Nothing on this sheet answers a graded question --- it is the syntax, so your time goes to the science. The worked versions live in the recitation decks and \texttt{bootcamp-notes.pdf}.}

\end{document}
